\subsection{CONTINUITY WITH RESPECT TO A SUBSPACE}

Let X and Y be two topological spaces, f a mapping of X into Y, B a subset of Y which contains \(f(\mathbf{X})\) . The definition of the induced topology as an initial topology (§ 2, no. 3, Proposition 4) shows that f is continuous at \(x \in X\) if and only if the mapping of X into the subspace B of Y, having the same graph as f, is continuous at x.

Now let A be a subset of X; if f is continuous at \(x \in A\) (resp. continuous on X), its restriction \(f|A\) is a mapping of the subspace A into Y, which is continuous at x (resp. continuous on A) by Proposition 2 of §2, no. 1. We shall sometimes say that a mapping \(f: X \to Y\) is continuous relative to A at \(x \in A\) (resp. continuous relative to A) if its restriction \(f|A\) is continuous at x (resp. continuous on A).

It should be noted that \(f|A\) can be continuous without \(f\) being continuous at any point of \(X\) ; an example of this phenomenon is provided by the characteristic function \(\varphi_A\) of a subset \(A\) of \(X\) which is such that both \(A\) and its complement are dense in \(X\) (§ 2, Exercise 11), \(\varphi_A\) being regarded as a mapping of \(X\) into the discrete space \(\{0, 1\}\) . \(\varphi_A\) is not continuous at any point of \(X\) , but its restriction to \(A\) is constant and therefore continuous.

If A is a neighbourhood in X of a point \(x \in A\) , and if \(f: X \to Y\) is such that \(f|A\) is continuous at x, then f is continuous at x; for each neighbourhood of x in A is a neighbourhood of x in X (local character of continuity).

PROPOSITION 4. Let \((\mathbf{A}_t)_{i\in I}\) be a family of subsets of a topological space \(\mathbf{X}\) whose interiors cover \(\mathbf{X}\) , or which is a locally finite closed covering of \(\mathbf{X}\) . Let \(f\) be a mapping of \(\mathbf{X}\) into a topological space \(\mathbf{X}'\) . If the restriction of \(f\) to each of the subspaces \(\mathbf{A}_t\) is continuous, then \(f\) is continuous.

For if \(\mathbf{F}'\) is a closed subset of \(\mathbf{X}'\) and if \(\mathbf{F} = \overline{f}^{1} (\mathbf{F}')\) , then \(\mathbf{F} \cap \mathbf{A}_i\) is closed in \(\mathbf{A}_i\) for each \(i \in I\) ( \(\S 2\) , no. 1, Theorem 1) and therefore \(\mathbf{F}\) is closed in \(\mathbf{X}\) by Proposition 3 of no. 1; the result now follows from Theorem 1 of \(\S 2\) , no. 1.

